\subsection{Threads are additional data}

The real zero tube has curve fibers.  To obtain moving points one must select them.

\begin{definition}[Finite zero thread]
\label{def:finite-zero-thread}
Let $\pi:\calZ\to B$ be an embedded real zero tube over a connected
one-manifold.  A \emph{finite zero thread} is a properly embedded
one-dimensional submanifold $K\subset\calZ$, with
$\partial K=K\cap\pi^{-1}(\partial B)$ when the base has boundary, such that
$\pi|_K:K\to B$ is a finite-sheeted covering.  For $B=S^1$, this forces $K$ to be
a compact one-manifold without boundary.  The pair $(\calZ,K)$ is a
\emph{threaded zero tube}.  A thread is \emph{intrinsic} only if its construction
from the stated AEG data is functorial under the declared notion of AEG
isomorphism.
\end{definition}

In the helical family, every slice $x=c$ gives
\[
 K_c=\calZ(A_{m,n})\cap\{x=c\},
\]
a degree-$2n$ finite thread.  Its homology and component structure are computed in
\eqref{eq:helical-boundary-class}.  The choice of $c$ is a probe choice; if one uses
$c=\pm L$, the boundary is part of the domain data but the choice of boundary
component remains.  Selecting only one of the $2d$ components is another choice.

A rank-two root incidence such as $\calR_F$ is already one-dimensional, so it plays
the role of its own finite thread.  If its moving points lie in a disc, it is an
ordinary geometric braid.  If they lie in an annulus or general surface, the
relevant surface braid group and equivalence relation must be named.

\subsection{A scalar Markov obstruction}

The first novelty filter is independent of AEG.

\begin{proposition}[Stateless additive collapse]
\label{prop:stateless-additive-collapse}
Let $A$ be an abelian group and $n\ge2$.  Every homomorphism
$\Theta_n:B_n\to A$ factors through the exponent sum (writhe)
$w:B_n\to\Z$:
\[
 \Theta_n(\beta)=w(\beta)c_n
\]
for a fixed $c_n\in A$.  If the family $\Theta_n$ is compatible with the standard
inclusions $B_n\hookrightarrow B_{n+1}$ and is unchanged under both Markov
stabilizations $\beta\mapsto\beta\sigma_n^{\pm1}$, then every $\Theta_n$ is zero.
\end{proposition}

\begin{proof}
In an abelian target, the braid relation
$\sigma_i\sigma_{i+1}\sigma_i=\sigma_{i+1}\sigma_i\sigma_{i+1}$ implies
$\Theta_n(\sigma_i)=\Theta_n(\sigma_{i+1})$.  All standard generators therefore
have the same image $c_n$, and the homomorphism is exponent sum times $c_n$.
Inclusion compatibility identifies $c_n$ with $c_{n+1}$.  Positive stabilization
then gives $c_{n+1}=0$; negative stabilization gives the same conclusion.  Hence
the family vanishes.
\end{proof}

Allowing a normalization depending on strand number does not evade both signs.  If
$I_n(\beta)=c\,w(\beta)+d n$ in an additive coefficient group, positive and
negative stabilization require $c+d=0$ and $-c+d=0$.  Away from two-torsion this
forces $c=d=0$, and in all cases it leaves no noncommutative information.  A useful
AEG invariant must retain state, use a trace-like normalization, or belong to a
more structured category.

\subsection{Affine conjugation and the Alexander layer}

Let $R$ be a commutative ring.  Write an affine transformation as $(p,u)$, with
$p\in R^\times$ and $u\in R$, acting by $x\mapsto px+u$.  Thus all multipliers
$p,q,r$ below are units and all translation coordinates $u,v,w$ lie in $R$.  We use
\begin{equation}
 (p,u)(q,v)=(pq,u+pv),
 \qquad
 (p,u)^{-1}=(p^{-1},-p^{-1}u).
 \label{eq:affine-group-law}
\end{equation}
Right conjugation $x\triangleright y=y^{-1}xy$ is
\begin{equation}
 (p,u)\triangleright(q,v)
 =\left(p,q^{-1}\bigl(u+(p-1)v\bigr)\right).
 \label{eq:affine-conjugation}
\end{equation}
On the fixed-multiplier branch
$X_t=\{(t,u):u\in R\}$ over a commutative ring in which $t$ is a unit, this becomes
\begin{equation}
 u\triangleright v=t^{-1}u+(1-t^{-1})v.
 \label{eq:alexander-quandle-operation}
\end{equation}
It is exactly an Alexander quandle, not a structure beyond the Alexander layer.

The affine torsion expression
\begin{equation}
 \kappa\bigl((p,u),(q,v)\bigr)=(1-q)u+(p-1)v
 \label{eq:affine-torsion}
\end{equation}
restricts to
\begin{equation}
 \kappa_t(u,v)=(t-1)(v-u).
 \label{eq:fixed-torsion}
\end{equation}

\begin{proposition}[Ordinary exactness]
\label{prop:ordinary-affine-exactness}
With the ordinary quandle coboundary convention
$df(x,y)=f(x\triangleright y)-f(x)$, one has
\[
 \kappa_t=d(tf)
 \qquad\text{for }f(u)=u.
\]
Thus the fixed-multiplier affine torsion represents the zero ordinary quandle
cohomology class.
\end{proposition}

\begin{proof}
Equation~\eqref{eq:alexander-quandle-operation} gives
\[
 df(u,v)=(t^{-1}-1)u+(1-t^{-1})v
 =\frac{t-1}{t}(v-u).
\]
Multiplication by $t$ gives \eqref{eq:fixed-torsion}.
\end{proof}

\subsection{Twisting, resonance, and a nonzero class}

We now use the twisted convention of Carter--Elhamdadi--Saito
\cite{CarterElhamdadiSaito2002TwistedQuandle}.  Let the coefficient action be a
scalar $T_{\mathrm{coef}}$.  A normalized twisted $2$-cocycle $\phi$ satisfies
\begin{align}
 T_{\mathrm{coef}}\phi(x,y)+\phi(x\triangleright y,z)
 ={}&T_{\mathrm{coef}}\phi(x,z)
 +(1-T_{\mathrm{coef}})\phi(y,z)\notag\\
 &+\phi(x\triangleright z,y\triangleright z),
 \label{eq:twisted-cocycle-equation}
\end{align}
and $\phi(x,x)=0$.  For coboundaries we choose the sign
\begin{equation}
 \delta_T^+f(x,y)=T_{\mathrm{coef}}f(x)
 +(1-T_{\mathrm{coef}})f(y)-f(x\triangleright y).
 \label{eq:twisted-coboundary-convention}
\end{equation}
The chain-level convention in the cited paper differs from
\eqref{eq:twisted-coboundary-convention} by a global minus sign in one displayed
definition, while its extension and state-sum calculations use the sign above.
The coboundary subgroup is the same under either choice.

Put $\alpha=t^{-1}$.  Direct substitution in
\eqref{eq:twisted-cocycle-equation} shows that $\kappa_t$ is a twisted cocycle for
every coefficient action.  If $\alpha-T_{\mathrm{coef}}$ is a unit in the
coefficient ring, it is exact: for
\[
 f(u)=\frac{t-1}{\alpha-T_{\mathrm{coef}}}u
\]
one has $\delta_T^+f=\kappa_t$.  Over a field, only the resonance
$T_{\mathrm{coef}}=\alpha$ survives this linear exactness test; over a general ring,
a nonzero nonunit difference may also obstruct this particular linear primitive.

\begin{theorem}[Finite-field resonant affine class]
\label{thm:resonant-affine-class}
Let $q$ be a prime power, let
$t\in\F_q^\times\setminus\{1\}$, and put $\alpha=t^{-1}$.  Give
$X_t=\F_q$ the Alexander operation
$x\triangleright y=\alpha x+(1-\alpha)y$ and give the coefficient module
$N=\F_q$ the resonant action $T_{\mathrm{coef}}=\alpha$.  Then
\begin{equation}
 \kappa_t(x,y)=(t-1)(y-x)
 \label{eq:finite-field-kappa}
\end{equation}
is a normalized twisted $2$-cocycle representing a nonzero class in
$H^2_{TQ}(X_t;N)$.
\end{theorem}

\begin{proof}
Normalization is immediate, and substituting the Alexander operation into
\eqref{eq:twisted-cocycle-equation} makes both sides equal
\[
 (t-1)\bigl[z-(\alpha+T_{\mathrm{coef}})x
 +(\alpha+T_{\mathrm{coef}}-1)y\bigr].
\]
Suppose $\kappa_t=\delta_T^+f$.  Subtracting the constant $f(0)$ does not change
the coboundary, so take $g(0)=0$.  Setting $y=0$ gives
\begin{equation}
 g(\alpha x)-\alpha g(x)=(t-1)x.
 \label{eq:resonant-recurrence}
\end{equation}
Let $r=\ord(\alpha)$.  Iterating \eqref{eq:resonant-recurrence} yields
\[
 g(\alpha^r x)=\alpha^r g(x)+r(t-1)\alpha^{r-1}x.
\]
But $\alpha^r=1$ and $r$ divides $q-1$, so the characteristic of $\F_q$ does not
divide $r$.  Taking $x\ne0$ contradicts $t\ne1$.  Hence no one-cochain has
coboundary $\kappa_t$.
\end{proof}

This closes the algebraic resonance question, but it does not yet produce a
stronger classical knot invariant.  In fact, a second calculation gives a sharp
collapse theorem.

We now fix the precise planar state-sum convention.  Give each oriented arc the
normal for which the ordered pair (tangent, normal) agrees with the plane
orientation, and give the unbounded region Alexander number zero.  At a crossing
$\tau$, let $L(\tau)$ be the Alexander number of the source region from which both
arc normals point.  Let $y_\tau$ be the over-arc color, and let its normal point
from an under-arc colored $x_\tau$ to one colored
$x_\tau\triangleright y_\tau$.  Define $\epsilon(\tau)=1$ when the ordered pair
(over-arc normal, under-arc normal) agrees with the plane orientation and
$\epsilon(\tau)=-1$ otherwise.  In additive coefficient notation the twisted
Boltzmann weight is
\begin{equation}
 B_T(\tau,C)=
 T_{\mathrm{coef}}^{-L(\tau)}
 \bigl(\epsilon(\tau)\kappa_t(x_\tau,y_\tau)\bigr)\in N,
 \label{eq:twisted-boltzmann-weight}
\end{equation}
and the state sum is
\begin{equation}
 \Phi_{\kappa_t}(D)
 =\sum_{C\in\operatorname{Col}_{X_t}(D)}
  \left[\sum_{\tau}B_T(\tau,C)\right]\in\Z[N].
 \label{eq:twisted-state-sum}
\end{equation}
Here $[a]$ denotes the basis element of the integral group ring indexed by
$a\in N$.  This is the additive form of the standard twisted convention.

\begin{theorem}[Planar state-sum collapse]
\label{thm:resonant-state-sum-collapse}
Under the hypotheses of Theorem~\ref{thm:resonant-affine-class}, the standard
twisted quandle cocycle state sum \eqref{eq:twisted-state-sum} is a Reidemeister
invariant of oriented classical links, but for every planar link diagram $D$ it is
\begin{equation}
 \Phi_{\kappa_t}(D)=
 \bigl|\operatorname{Col}_{X_t}(D)\bigr|[0].
 \label{eq:planar-state-sum-collapse}
\end{equation}
Thus it provides no enhancement of the Alexander quandle coloring count.
\end{theorem}

\begin{proof}
The Reidemeister invariance follows from the standard twisted-cocycle theorem
\cite{CarterElhamdadiSaito2002TwistedQuandle}.  To identify the value, define the
two-dimensional $\F_q$-vector space $E=N\oplus X_t$ with Laurent action
\begin{equation}
 T_E(a,x)=\bigl(\alpha a-(t-1)x,\alpha x\bigr).
 \label{eq:extension-action}
\end{equation}
This is invertible and makes
$0\to N\to E\to X_t\to0$ a short exact sequence of Alexander modules.  For the
section $s(x)=(0,x)$,
\begin{equation}
 T_Es(x)+(1-T_E)s(y)-s(x\triangleright y)
 =\bigl(\kappa_t(x,y),0\bigr).
 \label{eq:kappa-obstruction}
\end{equation}
Hence $\kappa_t$ is the obstruction cocycle of an Alexander-module extension.
The planar obstruction-cocycle cancellation theorem of
Carter--Elhamdadi--Saito then makes every coloring contribute the identity weight.
With convention \eqref{eq:twisted-state-sum}, the state sum is therefore exactly
$|\operatorname{Col}_{X_t}(D)|[0]$.
\end{proof}

Theorems~\ref{thm:resonant-affine-class} and
\ref{thm:resonant-state-sum-collapse} illustrate why cohomological nontriviality
and knot-theoretic detection must be tested separately.  A nonzero class can have a
trivial planar enhancement because its weights cancel through an extension.
On diagrams in nonplanar surfaces, the same extension argument need not cancel;
that is a different, surface-link question and is not promoted here to an ordinary
$S^3$ invariant.

\subsection{Variable multipliers and the Reidemeister-III anomaly}

For the full affine quandle, the torsion
\eqref{eq:affine-torsion} is generally not an ordinary $2$-cocycle.  With
$x=(p,u)$, $y=(q,v)$, and $z=(r,w)$, define its type-III defect by
\begin{align*}
 \mathfrak A_\kappa(x,y,z)
 ={}&\kappa(x,y)+\kappa(x\triangleright y,z)
 -\kappa(x,z)-\kappa(x\triangleright z,y\triangleright z).
\end{align*}
A direct calculation gives
\begin{equation}
 \mathfrak A_\kappa(x,y,z)
 =\frac{(q-r)(r-1)}{qr}\,\kappa(x,y).
 \label{eq:variable-multiplier-anomaly}
\end{equation}
It vanishes on a fixed-multiplier branch, as it must.  On a one-component knot,
conjugate meridians also have equal multipliers, so arbitrary multiplier variation
along arcs is not legitimate coloring data.  Nonzero use of
\eqref{eq:variable-multiplier-anomaly} would require additional region state, a
multi-component setting, a dynamical quandle, or another declared structure.

\begin{warning}[Anomaly is not associator]
A failure of the ordinary type-III cocycle equation is not automatically a higher
associator or a surface-holonomy invariant.  Such an interpretation requires
separate pentagon/hexagon or equivalent coherence and filling-independence proofs.
\end{warning}

\subsection{A historical free-word calculation}

One useful repository computation can be stated without overclaiming it.  Use the
historical figure-eight knot-group presentation
\[
 G_{4_1}=\langle a,b\mid abbbaBAAB=1\rangle,
 \qquad A=a^{-1},\quad B=b^{-1}.
\]
Under
the convention that $a:x\mapsto tx$, $b:x\mapsto x+1$, uppercase letters denote
inverses, and the displayed transformations compose in the fixed word order, the
free-group word
\[
 abbbaBAAB
\]
has unit multiplier and translation coordinate
\begin{equation}
 -\Delta_{4_1}(t)=-(t^2-3t+1).
 \label{eq:figure-eight-free-word}
\end{equation}
This is a verified affine word calculation.  The displayed generator assignment
descends from the free group to $G_{4_1}$ only on the parameter locus
$\Delta_{4_1}(t)=0$.  At generic $t$ it is not a knot invariant; it is the failure
of the chosen generator assignment to satisfy the relator.

\subsection{The remaining invariant gate}

Theorem~\ref{thm:aeg-braid-realization} supplies every braid, and
Definition~\ref{def:finite-zero-thread} supplies explicit threads in selected real
tubes.  Neither result selects a choice-free numerical or algebraic decoration.
For a family $I_n:B_n\to\mathcal A$ to descend to oriented links via Markov's
theorem, it must at minimum satisfy
\begin{equation}
\begin{aligned}
 I_n(\gamma\beta\gamma^{-1})&=I_n(\beta),\\
 I_{n+1}(\iota_n\beta\,\sigma_n)&=I_n(\beta),\\
 I_{n+1}(\iota_n\beta\,\sigma_n^{-1})&=I_n(\beta).
\end{aligned}
\label{eq:markov-gate}
\end{equation}
after any declared normalization \cite{Birman1974Braids}.  An AEG construction
must additionally prove independence of logarithmic lift, probe, thread, spine,
axis, basepoint, and expression presentation whenever those are auxiliary choices.

\begin{problem}[AEG Markov descent]
\label{prob:aeg-markov-descent}
Construct a nonabelian or stateful decoration natural under AEG isomorphisms, prove
the braid relations and all choice-independence statements, and find a Markov
normalization satisfying \eqref{eq:markov-gate}.
\end{problem}

\begin{problem}[Separation beyond Alexander and Burau]
\label{prob:beyond-alexander-burau}
After Problem~\ref{prob:aeg-markov-descent} is solved, exhibit two oriented links
with the same declared Alexander/Burau baseline but different AEG invariant.  A
stronger braid representation without closure descent does not meet this criterion.
\end{problem}

\subsection{Conclusion}

The singular-zero programme now has a linked analytic, arithmetic, and topological
core.  Holomorphic pullback turns degree-$m$ ramification into a $2m$-pronged
essential zero and a $2\pi m$ cone.  For the projective operators generated by
$z\mapsto z+\sqrt2$ and $z\mapsto-1/z$, the normalized Hecke Hauptmodul makes this
mechanism global: its $\{\infty,4\}$ dessin is an exact singular AEG zero network,
and its assignment descends with a nontrivial sign character.  This replaces a
postulated apeirogonal drawing by a zero graph derived from an arithmetic operator
group.

The same sign character has a precise but layered knot-theoretic meaning.  Its
coarse completed AEG carrier is the flat cylinder, whereas the unit tangent bundle
of the hyperbolic sign cover is, under the declared cusp compactification, the
complement of $T(2,4)$ in $S^3$.  Primitive hyperbolic operator classes give
classical Hecke periodic-orbit knots, with the zero dessin serving as their coding
spine only after the required directed thickening.  This identifies an exact
arithmetic--geometry--knot interface; it does not turn the local four-prong germ
into a link or produce a new knot invariant.

The polynomial threaded-carrier family adds the relation between slices that is
absent from a fixed zero dessin.  For $P_m(z,t)=z^2-t^m$, the real equation
$\im P_m=0$ sweeps a genus-$(m-1)$ surface through $2m$ saddle
reconnections, while $P_m=0$ is an intrinsic root thread closing to $T(2,m)$.
The discriminant order and logarithmic period, the two-braid exponent, and half
the negative Euler characteristic are the same integer $m$.  The exact sequence
$1\to P_2\to B_2\to S_2\to1$ makes the forgetting visible: at $m=4$ the
polynomial factors and its root permutation is trivial, yet the pure-braid
coordinate and the two-component linking number are both two.  The $q=4$ cusp
slope calculation selects this binomial as a peripheral toric normal form, while
the compactified sign cover recovers it more intrinsically as the cusp-divisor
section in its logarithmic-tangent section ring.  The corresponding weighted
Hopf circle bundle recovers the filled $T(2,4)$ cusp link.  This logarithmic cone
does not itself supply the radial genus-three carrier, and its arithmetic
quadratic twist depends on a declared descent datum.

The exact-sequence comparison also explains why the first two-strand attempt is
too small: its order-four relator has a nonprimitive pure-braid residue.  In four
strands, the Garside half twist squared and the four-point rotation raised to the
fourth power are the same full twist.  Hence the q=4 unit-tangent extension is a
pullback of the four-braid center extension, and the normalized discriminant
period recovers one fiber winding.  The resulting integral character on the
operator-kernel relations has a genuine Lyndon--Hochschild--Serre transgression:
its class is the unit-tangent Euler extension class.  This is not a transgression
over the circle base and is not Paper~I's affine torsion.  The four root paths
used in this calibration are supplied; deriving them from marked arithmetic
histories remains the naturality gate.

The relative-divisor theorem separates irreducibility over the coefficient field,
geometric sheets after scalar extension, and analytic monodromy away from the
discriminant.  Supplied quadratic and quartic registers now make that hierarchy
finite and computable: they give exact endpoint divisors, collapse kernels,
discriminants, monodromy, and Frobenius cycles, while a tagged trace retains some
intermediate domain data.  Their terminal divisor still factors through projective
evaluation.  Thus they are a restricted naturality test, not the general functor
from marked histories to prime relative divisors.

Downstream, regular real assignments produce curve tubes under properness; their
compact circle-parameter total spaces are tori.  The helical family makes component
transport, deck lifting, and selected threads explicit.  Singular AEG metrics
realize the two elementary real Morse events, while complex root fillings isolate
the branch half-twist.  Paper II's function theory realizes every braid by
arithmetic-holomorphic roots.

The same analysis locates the frontier.  Flexible metrics and the
$\sin(1/z)$ control make analytic existence abundant, while universal braid
realization makes mere representability nondiscriminating.
Stateless scalars collapse to writhe.  Fixed affine multipliers reduce to an
Alexander quandle; their ordinary torsion is exact.  Resonant twisting gives a
genuinely nonzero cohomology class, but its planar classical-link state sum still
collapses to the coloring count.  The remaining original problem is therefore not
to draw another tube or to relabel a classical periodic orbit.  It is to derive
canonical register or divisor data from marked histories, determine what survives
inside the operator-to-knot fibers, and prove that any resulting stateful
decoration survives every quotient required to become a knot invariant.
